\chapter{The Software Architecture}\label{ch:software}
\renewcommand{\thechapter}{A}

\lettrine{\color{MainColor}I}{n} parallel to the theoretical work, to make the experimental phase possible, a critical effort has been spent on the software definition and implementation for the real-time control architecture. The work started from the previous one developed during the TRIDENT project. However, its limitations and the need to use a new, recently developed real-time robotic middle-ware, have required a new implementation. 
Given these premises, the new architecture had the important goal to further improve the reuse of each software components. Toward that end, the methodology that I have followed is to subdivide the architecture in different processes, each one with a specific concern (e.g. Driver, Controller, Logger). 
The implementation of the RT framework has been done on top of the robotic middle-ware developed at GRAAL (the Ortos RT Framework). However, to facilitate the interaction with other software components that do not depend on real-time requirements, a ROS interface was also built to act as a bridge between the RT and non-RT environments. 
The software includes single and bimanual arm control, coordinated object transportation, leader/follower path following, all running through the task priority control module. The above-mentioned framework is structured in a series of dedicated processes, communicating via internal API or through a ROS interface. Having to deal with two task schedulers at the same time (Ortos and ROS), the integration of the two has not been so trivial, and many hours where spent in fighting with segmentation faults and broken processes, but that's a recurring old story. The architecture can be summarized as shown in \ref{fig:softwareArch}: each box corresponds to a process, following a modular design to ensure re-usability and real-time performances dependability.

\begin{figure}
	\centering
	\includegraphics[width=0.9\textwidth]{softwareArch/software_architecture.eps}
	\caption{A simplified overview of the developed software where three central modules of the control architecture have been highlighted in blue: the Controller, the Driver and the Logger.}
	\label{fig:softwareArch}
\end{figure}


\section{Controller Modules}
As above mentioned the software architecture is composed by different modules that run as soft real time tasks of the Ortos RT middle-ware scheduled in a round-robin fashion at a $\SI{100}{\hertz}$ frequency rate. The Ortos middle-ware provides all the utilities to implement message queues and shared memory read/write operations. In figure \ref{fig:softwareArch} it is showed a simplified overview of the overall architecture, with some of its central modules being highlighted. Let's now see a little bit more in details how they operate. The nine tasks that have been developed for creating the software architecture are the following:

\begin{itemize}[label=\termItemize]%
	%\setlength\itemsep{-.05em}%
	\item \emph{KCL Controller}: the hearth of the software architecture, where all the heavy maths are going on.
	\item \emph{Driver}: the direct link between the software and the robot's hardware
	\item \emph{Logger}: listens for all the data-exchanged among the modules and stores all the variable of interest
	\item \emph{Data Monitor}: a collection of all the relevant system variable to check the system status at a glance
	\item \emph{Command Console}: a command line interface to generate simple commands
	\item \emph{Keyboard Controller}: an interface to generate references for the controller in a tele-operated fashion
	\item \emph{UDP Receiver}: the module where data can be received from other computers running the software architecture 
	\item \emph{UDP Sender}: the module where data can be sent to other computers running the software architecture 
	\item \emph{MoCap Bridge}: the link between the motion capture system and the control architecture
\end{itemize}

Let's now spend just few words going through these nine modules and seeing how they relate one another.

\subsubsection{KCL Controller}
This first process is the principal control module, where all the information about the actual state of the system is gathered, and the task priority inverse kinematics takes place. It is the end-point for the controller interface, that has been implemented in the \texttt{ControllerInterface} class. This class exposes all the functionalities of the framework for an easy integration of the KCL with the HLL, thereby enforcing the \ref{c:interfaces} contribution.
This module loads the model of the robot at hand and stores it in the \texttt{ArmModel} class. The model smart pointer is passed through classes so to always have an updated status of the model available across the whole control architecture. It communicates mainly with the \textit{Driver} from which it receives the feedback and other state relevant variables such as: enable status of the robot, availability of grippers, status of the calibration and more. After all the model loading and other bootstrap procedures, the task  When a command has been issued by the control interface, once the control has been computed (via the \texttt{MultiArm} class), it sends the control reference back to the \textit{Driver}.
This module also provides an alive signal which tells whether the process is actually running and if the initial bootstrap has been successful. This information is used by other modules to provide relevant warnings and also in the graphical user interface (GUI) which is presented in section \ref{sec:gui}.

\subsubsection{Driver}
The driver is the module which directly communicates with the robot and exposes an interface which makes the control of the robot transparent to all the other modules, the \textit{Controller} in particular. It can have internal different implementations depending on the way we access the robot actuators and sensor. Let's take some examples based on the Baxter and YouBot platform. On the Baxter we have no access to the robot internal hardware due to proprietary restrictions, and this modules is implemented as a ROS wrapper. In the YouBot case instead we had access to the hardware APIs, and in the driver all the base and gripper calibration had to be implemented. The \textit{Driver} is mainly composed of the \texttt{Robot::Driver} class which is subdivided in the \texttt{Robot::Base} and \texttt{Robot::Arm} classes.
Also this module provides an alive signal which tells whether the process is actually running and if the initial robot procedures have been successful.

\subsubsection{Logger}
The logging module is one the most important modules, as without it our experiment would just be overly complicated diversions. This task in particular is the only for which a slower scheduling rate has been chosen, due to the high input/output operations throughput. This was particularly the case for the Baxter case where, having two 7 DOFs arms we have a massive amount of data to be saved each iteration. All the data is read from the Ortos shared memory and includes: robot parameters, robot feedback (position, velocity, control mode), references evaluated by the controller, task priority relevant variables for each task (intermediate rhos, activations, enable).
Using the console it is possible to start and stop the logging.

\subsubsection{Data Monitor}
The \textit{Data Monitor} shows a collection of all the relevant system variables to check the system status at a glance, reading all the data from the Ortos shared memory written by the other tasks. It has been of great help during the development of the code to check what was going on when the system was acting like crazy for some insidious error, which as you would expect happened 80\% of the time.

\subsubsection{Command Console}
The \textit{Command Console} is command line interface to generate simple commands to control the robot. It can be very useful when testing single functionalities of the framework during implementation, but can also serve as a control interface for simple tasks. Via this modules you can access most of the capabilities offered by the controller. It has been developed both as interactive shell and as a terminal CLI (to provide command history). Just as an example, the typical screen for the command console can be the following:

\begin{lstlisting}[style=customc]
jointPos             < arm j1 ... jn >
jointVel             < arm j1 ... jn >
cartPos              < arm y p r x y z >
cartRelToolPos       < arm y p r x y z >
cartVel              < arm y p r x y z >
cartVelInput         < arm >
stopArm              < arm >
openGrip             < arm >
closeGrip            < arm >
reloadPar            < arm >
holdMode             < arm [0(disable)|1(enable)] >
startLog             < arm >
stopLog              < arm >
unfoldArm            < arm >
foldArm              < arm >

> Type a command...
\end{lstlisting}

In the interactive shell version of the console also, all the messages issued by the other modules are displayed in a unified way, clearly identified with a sender name and colour, to have a single place where to look for knowing the system's state.

\subsubsection{UDP Sender and Receiver}
The \textit{UDP Sender} module, together with \textit{UDP Receiver} one, became necessary for the realization of the YouBots experiments, where multiple PCs had to communicate one another, as summarized earlier in Figure \ref{fig:youbotarchitecture}. Depending on a series of input parameters these module can be configured to run for reading agent related inputs (as on the YouBots) or external coordination data (as in the Central Server of the above mentioned figure). The receiver module \bsq{simply} opens a series of non-blocking UDP sockets and writes the received data to the Ortos framework. The sender act in a complementary way.

\subsubsection{Keyboard Controller}
This task is a simple keyboard controller that can be used to generate references for the KCL in a tele-operated fashion. It makes use of the SDL Library, a well known library for input management, that gives the possibility to control the robot in a WASD-like mode. Via this controller is possible to switch between sending smoothly increasing and decreasing references (implemented with a simple first order filter), or step-wise references. 


\subsubsection{MoCap Bridge}
The motion capture task serves as a bridge between the OptiTrack{\scriptsize \raisebox{.3em}{\texttrademark}} system and the control framework. In particular, it reads the rigid body data streamed by the MoCap proprietary software and translates it into coordinate frames compatible with the mathematical library. The data is then sent to the \textit{UPD Receiver} tasks of each agent. One of the main issues has also been the correct axis interpretation, as everyone uses its own \textit{standard}, and getting it right isn't always an easy task.


\paragraph{}
Let's note how in the last YouBot experiments for example, where we have multiple PCs composing the overall system, the architecture running is the same for all the sub-systems. Of course the controller and driver modules will be only running on the agents, but the general software structure is shared, with just different configuration parameters being parsed. Also  the definition of a standardized API interface for such processes allows the architecture to be as much as possible generic and reusable for different setups.

%\section{Classes and relationships}
%
%\begin{figure}[H]
%	\begin{center}
%		%\leavevmode
%		\includegraphics[width=.6\textwidth]{softwareArch/classCTRL_1_1MultiArm__coll__graph}
%		\caption{The collaboration diagram of the \texttt{MultiArm} class, generated by the \textit{Doxygen} documentation tool.  }
%	\end{center}
%\end{figure}
%
%\begin{itemize}
%	\setlength\itemsep{-.3em}
%	\item \texttt{ControllerInterface}
%	\item \texttt{SingleArm}
%	\item \texttt{MultiArm}
%	\item \texttt{TaskPriority}
%	\item \texttt{ArmModel}
%	\item \texttt{VehicleModel}
%	\item \texttt{JacobianHelper}
%\end{itemize}

\subsection{External Libraries}\label{sec:libs}
To make the controller development possible, a series of external libraries as been included. These are all the libs that have been used:

\begin{itemize}
	\setlength\itemsep{-.3em}%
	\item {\bfseries Ortos} framework - Realtime middleware
	\item {\bfseries ROS Indigo} - The ROS Operating system
	\item {\bfseries CMAT} - Robotics mathematical library
	\item {\bfseries Libconfig 1.5} - Configuration file manager
	\item {\bfseries SDL 1.2} - Media library for user input
	\item {\bfseries Beep} - Utility program for acoustic signalling
\end{itemize}

To improve the usability of the software, a GUI interface has been developed to interact with the controller whose details are included in the appendix.  

\section{Graphical User Interface}\label{sec:gui}
During the PhD work I also developed a simple user interface. The reasons that lead to this decision are fundamentally three: a personal interest in interface design, speed up the testing phase without having to hassle with too many open terminals (and whoever works in applied computer science knows what I mean), and last but not least providing a tool that was easy to use and fool's proof for whoever had the need of using such framework. The interface was developed in QT5 using the new QML scripting language. QML is basically a JavaScript-like language for developing cross-platform window applications and Model-View-Controller relationships.

It all started with a simple \framebox[1.08\width]{\textsc{Emergency-Stop}} button to stop the robot in case of critical errors, and lately I expanded the graphical interface to encompass the following features:

\begin{itemize}
	\setlength\itemsep{-.3em}
	\item Enable/Disable the robot
	\item Provide status indicators
	\item Sending control commands (with type checking)
	\item View Command History
	\item View robot feedbacks
	\item Plot real-time graphs (still in beta phase)
\end{itemize}

The interfaces where developed for the Baxter and the Youbot robots. They share the same code base, apart from some robot dependent customization. In Figure \ref{fig:guibaxter} are showed some self-explanatory screenshots of the interface for the Baxter Robot.

\begin{figure}[]
	\centering
	\subcaptionbox{\label{fig:guibaxt1}}{\includegraphics[width=0.48\columnwidth]{softwareArch/GUI_1.png}}
	\hspace{0.25cm}
	\subcaptionbox{\label{fig:guibaxt2}}{\includegraphics[width=0.48\columnwidth]{softwareArch/GUI_3.png}}\\
	\subcaptionbox{\label{fig:guibaxt3}}{\includegraphics[width=0.48\columnwidth]{softwareArch/GUI_4.png}}
	\hspace{0.25cm}
	\subcaptionbox{\label{fig:guibaxt4}}{\includegraphics[width=0.48\columnwidth]{softwareArch/GUI_5.png}} \\
	\subcaptionbox{\label{fig:guibaxt5}}{\includegraphics[width=0.48\columnwidth]{softwareArch/GUI_6.png}}
	\hspace{0.25cm}
	\subcaptionbox{\label{fig:guibaxt6}}{\includegraphics[width=0.48\columnwidth]{softwareArch/GUI_7.png}}
	\caption{Various states and tabs of the GUI.} 
	\label{fig:guibaxter}
\end{figure}

Let's briefly describe the various states of the interfaces showed in Figure \ref{fig:guibaxter}. The interface is composed of three main tabs: Commander, Feedbacks ad Graphs. The original Emergency-Stop button is very wide and always accessible at the bottom of the interface, so to be functional even in case of a critical situation. Together with the stop button at the bottom of the interface we find the system status indicators for: the controller, the driver and the robot. In figure \ref{fig:guibaxt1} we can see the interfaces when no active controller is detected in the system, note all the red indicators in the bottom.
In the other figure instead the Driver and Controller have been activated, the status indicators became green and the interface becomes active. Figure \ref{fig:guibaxt2}  shows the available commands to be sent. The interface is equipped with a syntax spelling checker which informs the user if a sent command was meaningful or not, informing it with a pop-up message. In figure \ref{fig:guibaxt3}  we can see the message resulting from a correct parsing of the command, whose syntax is available with the button \framebox[1.1\width]{\textsc{View Syntax}} on the right (as seen in \ref{fig:guibaxt4}). Figure \ref{fig:guibaxt5}  shows the feedback tab. Figure \ref{fig:guibaxt6}  instead shows the prototype tab for plotting related features. 


